immini[thm]{\subsubsection{Quá trình đẳng áp}
\newghichu{\textit{Quá trình đẳng áp} là quá trình biến đổi trạng thái của chất khí trong đó áp suất được giữ nguyên không đổi}
\subsubsection{Định luật Charles}
\newghichu{Khi áp suất của một khối lượng khí xác định giữ không đổi thì thể tích của khí tỉ lệ thuận với nhiệt độ tuyệt đối của nó.
\begin{align*}^{}
	V\sim T\ hay\ \dfrac{V}{T}=\ct{ hằng số}
\end{align*}}
}{includegraphics[width=4cm]{Images/charleslaw}}

\Binhluan{
\begin{itemize}
	\item[] Khi chất khí chuyển từ trạng thái này sang trạng thái khác thì ta có:
\begin{align*}\boder{\dfrac{V_1}{T_1}=\dfrac{V_2}{T_2}= \dfrac{V_3}{T_3}=..}\end{align*}
\end{itemize}

}{Thể tích V của khí không tăng tỉ lệ thuận với nhiệt độ Celsius}
\subsubsection{Đường đẳng áp}
\newghichu{Đường đẳng áp là đường biểu diễn sự phụ thuộc của thể tích theo nhiệt độ khi áp suất không đổi}
\begin{itemize}
	\item Dạng đường đẳng áp trong các hệ trục toạ độ:
	\begin{center}
		\begin{tikzpicture}
			\draw[->] (0,0)node[left]{O}--(3,0)node[below]{$T$}; 
			\draw[->] (0,0)--(0,3)node[left]{V};
			\draw[dashed] (0,0)--++(50:1);
			\draw [blue, thick](0,0)++(50:1)--++(50:2)node[above,black]{$p_1$};
			\draw[dashed] (0,0)--++(30:1);
			\draw[blue, thick] (0,0)++(30:1)--++(30:2)node[right,black]{$p_2$};
			\draw (0,0)++(40:3)node[above right]{$p_1<p_2$};
			\node at (1.5,-1){Hệ trục toạ độ (V, T)};
			\begin{scope}[shift={(5,0)}]
				\draw[->] (0,0)node[left]{O}--(3,0)node[below]{$p$}; 
				\draw[->] (0,0)--(0,3)node[left]{$V$};
				\draw [blue, thick] (2,1)--++(90:2);
				\node at (1.5,-1){Hệ trục toạ độ (p, V)};
			\end{scope}
			\begin{scope}[shift={(10,0)}]
				\draw[->] (0,0)node[left]{O}--(3,0)node[below]{$p$};
				\draw[->] (0,0)--(0,3)node[left]{$T$};
				\draw [blue, thick] (2,1)--++(90:2);
				\node at (1.5,-1){Hệ trục toạ độ (T, p)};
			\end{scope}
		\end{tikzpicture}	
	\end{center}
	\begin{note}
		Trong hệ tọa độ (V, T) đường đẳng áp là đường thẳng kéo dài đi qua gốc tọa độ.
	\end{note}
\end{itemize}
\subsubsection{Các định luật Boyle và Charles là các định luật gần đúng}
\begin{itemize}
	\item Các định luật Boyle và Charles được rút ra từ những thí nghiệm thực hiện trong điều kiện áp suất không quá $10^6\ Pa$, nhiệt độ không dưới 200 K. 
	\item Các thí nghiệm thực hiện trong điều kiện áp suất rất cao và nhiệt độ rất thấp cho kết quả không phù hợp với các định luật trên.
	\item Để phân biệt khí lí tưởng với khí thực người ta định nghĩa khí lí tưởng là khí tuân theo đúng các định luật Boyle và Charles.
	\item Tuy nhiên, sự khác biệt giữa khí lí tưởng và khí thực không lớn ở điều kiện bình thường về áp suất và nhiệt độ nên người ta vẫn có thể áp dụng các định luật của khí lí tưởng cho khí thực nếu không cần độ chính xác cao.
\end{itemize}